\scanpage{61}
% ll. 1–4
Notons que, si $\unsure{Z_i|z_i}$ \inserted{$\subset\operatorname{GL}(V)$} est le groupe engendré par $\tau_i$, on a $Z_i\cap k^\times=1$ donc $Z_i\hookrightarrow\operatorname{Aut}(C)$, donc $\nu_i$ prend une des valeurs $2,3,4,5$
\[
2\leq\nu_0,\nu_1\leq5
\]

% ll. 5–8
\deleted{Il faut} Plus gén., \deleted{les valeurs propres des \unsure{$\tau_i$ engendrant}} \inserted{pour tte ps. réfl. \inserted{stricte} de $G_{0}$, son} \deleted{\unsure{partic.} des groupes engendrés par} ordre $\nu$ est comprise entre 2 et 5 \ldots
% archive sheet number: 60
